\documentclass[10pt,a4paper]{amsart}
\usepackage{amsmath,amssymb,amsthm,graphicx}
\usepackage{verbatim}
\newenvironment{reference}{\comment}{\endcomment}
%\newenvironment{reference}{}{}
\newenvironment{slogan}{\comment}{\endcomment}
\newenvironment{history}{\comment}{\endcomment}

% For commutative diagrams we use Xy-pic
\usepackage[all]{xy}

% We use 2cell for 2-commutative diagrams.
\xyoption{2cell}
% Theorem environments.
%
\theoremstyle{plain}
\newtheorem{theorem}[subsection]{Theorem}
\newtheorem{proposition}[subsection]{Proposition}
\newtheorem{lemma}[subsection]{Lemma}

\theoremstyle{definition}
\newtheorem{definition}[subsection]{Definition}
\newtheorem{example}[subsection]{Example}
\newtheorem{exercise}[subsection]{Exercise}
\newtheorem{situation}[subsection]{Situation}

\theoremstyle{remark}
\newtheorem{remark}[subsection]{Remark}
\newtheorem{remarks}[subsection]{Remarks}

\numberwithin{equation}{subsection}

% Macros
%
\def\lim{\mathop{\mathrm{lim}}\nolimits}
\def\colim{\mathop{\mathrm{colim}}\nolimits}
\def\Spec{\mathop{\mathrm{Spec}}}
\def\Hom{\mathop{\mathrm{Hom}}\nolimits}
\def\Ext{\mathop{\mathrm{Ext}}\nolimits}
\def\SheafHom{\mathop{\mathcal{H}\!\mathit{om}}\nolimits}
\def\SheafExt{\mathop{\mathcal{E}\!\mathit{xt}}\nolimits}
\def\Sch{\mathit{Sch}}
\def\Mor{\mathop{\mathrm{Mor}}\nolimits}
\def\Ob{\mathop{\mathrm{Ob}}\nolimits}
\def\Sh{\mathop{\mathit{Sh}}\nolimits}
\def\NL{\mathop{N\!L}\nolimits}
\def\CH{\mathop{\mathrm{CH}}\nolimits}
\def\proetale{{pro\text{-}\acute{e}tale}}
\def\etale{{\acute{e}tale}}
\def\QCoh{\mathit{QCoh}}
\def\Ker{\mathop{\mathrm{Ker}}}
\def\Im{\mathop{\mathrm{Im}}}
\def\Coker{\mathop{\mathrm{Coker}}}
\def\Coim{\mathop{\mathrm{Coim}}}

% Boxtimes
%
\DeclareMathSymbol{\boxtimes}{\mathbin}{AMSa}{"02}

%
% Macros for moduli stacks/spaces
%
\def\QCohstack{\mathcal{QC}\!\mathit{oh}}
\def\Cohstack{\mathcal{C}\!\mathit{oh}}
\def\Spacesstack{\mathcal{S}\!\mathit{paces}}
\def\Quotfunctor{\mathrm{Quot}}
\def\Hilbfunctor{\mathrm{Hilb}}
\def\Curvesstack{\mathcal{C}\!\mathit{urves}}
\def\Polarizedstack{\mathcal{P}\!\mathit{olarized}}
\def\Complexesstack{\mathcal{C}\!\mathit{omplexes}}
% \Pic is the operator that assigns to X its picard group, usage \Pic(X)
% \Picardstack_{X/B} denotes the Picard stack of X over B
% \Picardfunctor_{X/B} denotes the Picard functor of X over B
\def\Pic{\mathop{\mathrm{Pic}}\nolimits}
\def\Picardstack{\mathcal{P}\!\mathit{ic}}
\def\Picardfunctor{\mathrm{Pic}}
\def\Deformationcategory{\mathcal{D}\!\mathit{ef}}
\usepackage{cleveref}
\crefname{originallemma}{original lemma}{original lemmas}
\Crefname{originallemma}{Original lemma}{Original lemmas}
\makeatletter
\newlabel{lemma-NAK}{{00DV}{source}{}{sourcepointer-1}{}}
\newlabel{lemma-NAK@cref}{{[originallemma][0][]00DV}{[1][0][]source}{}{}{}}
\makeatother
\crefname{originallemma}{original lemma}{original lemmas}
\Crefname{originallemma}{Original lemma}{Original lemmas}
\makeatletter
\newlabel{lemma-permute-xi}{{00LJ}{source}{}{sourcepointer-2}{}}
\newlabel{lemma-permute-xi@cref}{{[originallemma][0][]00LJ}{[1][0][]source}{}{}{}}
\makeatother
\title{A good element is a nonzerodivisor and its quotient is Cohen-Macaulay with a shorter maximal regular sequence}
\author{The Stacks Project Authors}
\date{stacks-project:algebra.tex:local-snapshot-20261012; GNU FDL 1.2 or later}
\begin{document}
\maketitle
\noindent\textit{Editorial note.} All delivered mathematical text is the author's own, copied from the pinned Stacks Project chapter file algebra.tex; no statement or proof step was rewritten, reordered or removed.\par\medskip
\noindent\textit{Editorial note.} Preamble provenance, as required for this source: the frozen snapshot stages the chapter's own preamble file beside it (preamble.tex, 217 lines, 8116 bytes, the same shared Stacks-wide preamble the published chow and etale-cohomology cards use); the first line of the pinned chapter file is \textbackslash{}input\{preamble\}, so no preamble text lives inside the chapter itself. The card takes that file's lines 8-18 (the verbatim package and the three \textbackslash{}comment-based environments reference/slogan/history at lines 9-12, section 13.2, together with the Xy-pic setup) and its lines 154 to the runtime end of the file (the native theorem declarations and mathematical macros). Not one character was changed, no package was added, no macro was defined, and the excerpt serves only as preamble, never as a span coordinate.\par\medskip
\noindent\textit{Editorial note.} The same-chapter pointers are located at their \textbackslash{}label lines inside the pinned chapter file; the cross-chapter pointers are located in the project's own tag list staged in the snapshot (sources/stacks-chow/tags.txt), at the verbatim lines carrying their canonical tags.\par\medskip
\noindent\textit{Editorial note.} The counted proof refers to two same-chapter lemmas by label (Nakayama and the permutation of a regular sequence); both are recorded as navigation stubs located at their \textbackslash{}label lines, and the arguments consuming them stay inside the counted span.\par\medskip
\noindent\textit{Editorial note.} The counted spans contain no \textbackslash{}cite, so the card carries no bibliography.\par\medskip
\noindent\textit{Allowed standard prerequisites (editorial).}\begin{itemize}
\item Nakayama's lemma. Invoked in the counted proof at line 25218. Turns the equality f\_1 K = K for a finite-length kernel K into K = 0.
\item Permuting an M-regular sequence. Invoked in the counted proof at line 25228. Upgrades the regularity on the quotient to regularity of the pair g, f\_1 on M.
\end{itemize}\medskip
\ifdefined\hypertarget\hypertarget{sourcepointer-1}{}\fi
\noindent\textit{Source reference.} Original lemma 00DV: sources/stacks-algebra/pinned-original-source.tex, source lines 3833--3833. Navigation pointer for \textbackslash{}ref\{lemma-NAK\} used inside the counted span (line 25218); the argument that consumes it is retained in the span.\par\medskip
\ifdefined\hypertarget\hypertarget{sourcepointer-2}{}\fi
\noindent\textit{Source reference.} Original lemma 00LJ: sources/stacks-algebra/pinned-original-source.tex, source lines 16856--16856. Navigation pointer for \textbackslash{}ref\{lemma-permute-xi\} used inside the counted span (line 25228); the argument that consumes it is retained in the span.\par\medskip
\begin{lemma}
\label{lemma-good-element}
Notation and assumptions as above. If $g$ is good with respect to
$(M, f_1, \ldots, f_d)$, then (a) $g$ is a nonzerodivisor on $M$,
and (b) $M/gM$ is Cohen-Macaulay with maximal regular
sequence $f_1, \ldots, f_{d - 1}$.
\end{lemma}

\begin{proof}
We prove the lemma by induction on $d$.
If $d = 0$, then $M$ is finite and there is no case
to which the lemma applies.
If $d = 1$, then we have to show that $g : M \to M$ is
injective. The kernel $K$ has support $\{\mathfrak m\}$
because by assumption $\dim \text{Supp}(M) \cap V(g) = 0$.
Hence $K$ has finite length. Hence $f_1 : K \to K$ injective
implies the length of the image is the length of $K$, and hence
$f_1 K = K$, which by Nakayama's Lemma \ref{lemma-NAK} implies $K = 0$.
Also, $\dim \text{Supp}(M/gM) = 0$ and so $M/gM$ is Cohen-Macaulay
of depth $0$.

\medskip\noindent
Assume $d > 1$. Observe that $g$ is good for $(M/f_1M, f_2, \ldots, f_d)$,
as is easily seen from the definition. By induction, we have that
(a) $g$ is a nonzerodivisor on $M/f_1M$ and
(b) $M/(g, f_1)M$ is Cohen-Macaulay with maximal regular sequence
$f_2, \ldots, f_{d - 1}$. By
Lemma \ref{lemma-permute-xi}
we see that $g, f_1$ is an $M$-regular sequence.
Hence $g$ is a nonzerodivisor on $M$ and
$f_1, \ldots, f_{d - 1}$ is an $M/gM$-regular sequence.
\end{proof}

\end{document}
